\chapter{Conception et modelisation}

\section{Introduction}
Ce chapitre decrit la conception de la solution, la structuration des donnees et les principaux flux fonctionnels.

\section{Conception fonctionnelle}
La solution s'articule autour de deux parcours principaux :
\begin{itemize}
  \item parcours client (consultation, profilage, chat intelligent) ;
  \item parcours administrateur (authentification, gestion contenus, suivi analytics).
\end{itemize}

\section{Modelisation de la base de donnees}
La base PostgreSQL est concue pour centraliser les informations hotelieres et les traces d'utilisation :
\begin{itemize}
  \item \texttt{hotels}, \texttt{hotel\_settings} ;
  \item \texttt{conversations}, \texttt{messages} ;
  \item \texttt{guest\_profiles}, \texttt{analytics\_events} ;
  \item \texttt{reservations}, \texttt{events}, \texttt{attractions} ;
  \item \texttt{admin\_users}.
\end{itemize}

\section{Conception des interfaces}
Les ecrans principaux incluent :
\begin{itemize}
  \item page d'accueil multi-hotels ;
  \item page detail hotel ;
  \item interface de chat temps reel ;
  \item tableau de bord administrateur ;
  \item interface analytics.
\end{itemize}

\section{Flux de donnees}
Le flux d'un message suit les etapes suivantes :
\begin{enumerate}
  \item reception et validation de la requete ;
  \item enrichment via contexte hotel + historique session ;
  \item generation de reponse par le modele IA ;
  \item persistence des traces et evenements analytics.
\end{enumerate}

\section{Conclusion}
Cette conception fournit une base solide pour la mise en oeuvre technique. Le chapitre suivant detaille l'implementation du chatbot intelligent et son integration RAG.
